\section{Introduction}
Model-driven software engineering uses metamodels to define the structures of
modeling languages. Essential MOF (EMOF) supplies a shared vocabulary of classes,
properties, inheritance and associations. A reusable formal foundation should
explain which declarations are legal, which finite object structures they admit,
and which consequences follow for every supported metamodel. Such an account can
support later reasoning about model transformations, reflective operations and
interchange without rebuilding structural conformance for each language.

The difficulty is partly interpretive. EMOF combines its own restrictions with
adopted UML definitions and reflective-service clauses. A class-creation
prerequisite need not be a global metamodel restriction. Repeated references need
not denote distinct containers. Inherited properties must retain declaration
identity across different paths. A mechanized checker can be correct for a
predicate that resolves these questions incorrectly; source justification and
proof of implementation correctness therefore require separate evidence.

Formal MOF accounts already provide reusable results.
Boronat and Meseguer give an executable algebraic account of model types and
conformance \cite{boronat_meseguer_2010}. Coq4MDE proves model-composition
preservation under explicit hypotheses \cite{kezadri_etal_2014}; mechanized
UML/OCL and textual metamodeling also provide substantial precedents
\cite{brucker_tuong_wolff_2020,jouault_bezivin_2006}.
Our contribution is a version-specific, inspectable structural EMOF account
connecting source interpretations, occurrence-sensitive consequences and a
mechanized realization. We do not claim the first formal MOF semantics, and
choice of proof assistant or concrete notation alone is not the contribution.

We investigate three research questions:
\begin{description}
\item[RQ1.] Which structural interpretation is justified by the EMOF source set,
and which contextual distinctions change admitted schemas or snapshots?
\item[RQ2.] What reusable consequences and executable/authoring correspondence
can be established for that interpretation?
\item[RQ3.] What modeling consequences arise in a published language fragment,
and how do outcomes and execution costs compare with aligned EMF validation?
\end{description}

VL-MOF contributes (i) an explicit structural interpretation of EMOF 2.5.1,
with source-derived separating examples; (ii) a Lean library of conformance
predicates and consequences, with checker correctness and conditional
source-elaboration adequacy; and (iii) a realization evaluated through a
published Train metamodel projection and a declared EMF comparison protocol.
The semantic domain includes finite nested packages, classes and inheritance,
properties and binary associations, selected primitive/enumeration values,
occurrence multiplicity, uniqueness, ordering, opposites and static containment.
Defaults and unset history, reflective execution, richer values and full
self-description are deferred. Structural acceptance does not establish a
language's domain invariants or the correctness of its generated software.

Section~2 motivates these distinctions with the published case. Section~3
explains source interpretation, definitions and proofs; Section~4 describes
the realization. Section~5 answers the research questions from the evidence.
Section~6 discusses related work and threats to validity, and Section~7 identifies
the results and concrete extensions of this foundation.
